\chapter{\rm\bfseries Conclusion and Future Work}
\label{ch:conclusion}

\section{Answers to the Research Questions}

This thesis presented a controlled benchmark of six gripper variants on a
single platform (Franka FR3), measuring stationary contact force and grasp
outcome as a function of press depth and tilt angle over 845 analysed
trials. Returning to the four research questions posed in
Chapter~\ref{ch:introduction}:

\begin{enumerate}
  \item \emph{How does contact force vary with press depth for each finger
        design?} Every finray variant follows a single shared
        rise--peak--valley--recovery curve (pairwise NMSE below $1\%$
        over the full 10\,mm calibrated window), whereas the rigid
        parallel jaws ramp steeply toward the platform's excess-force
        limit within roughly 2\,mm of first contact.
  \item \emph{How much can the gripper be tilted before the grasp becomes
        unreliable?} The finrays grasp reliably across every orientation
        tested ($\pm 2.5^{\circ}$ about both axes, plus a $5^{\circ}$
        seat mount) over the full $\geq 10$\,mm window; the rigid jaws
        never exceed 2.5\,mm, and two stock-jaw tilt conditions produced
        no successful grasp at any depth.
  \item \emph{Does tilt tolerance depend on press depth?} Yes ---
        robustness is depth-gated: near the contact boundary, tilt
        strongly modulates both force and success (including a
        $2.4\times$ direction-dependent asymmetry about the geometrically
        asymmetric axis), whereas past the buckling threshold the mean
        across-orientation force spread shrinks to $1.3$--$2.4$\,N.
  \item \emph{What is the compliant finger's response under compression?}
        Blade buckling: force peaks at buckling onset ($\approx 2$--$3$\,mm
        past first successful grasp), dips as the blades buckle, then
        recovers --- and the finray geometry (blade thickness, seat
        angle, print instance) sets the force \emph{scale} (30--51\,N
        peaks) but not the curve \emph{shape}.
\end{enumerate}

The depth-dependent tilt robustness identified here had not been
quantitatively characterised in prior literature, to the author's
knowledge, and constitutes the principal empirical contribution.

\section{Implications}

With the robot, control policy, cloth, and environment held identical,
a finger swap alone buys at least a $4\times$ wider press-depth tolerance than the best rigid
configuration ($\geq 10$\,mm versus 2.5\,mm on the padded jaws, and
roughly $7\times$ versus the stock jaws), at roughly half the grip force.
Under tilt the advantage widens further, to $20\times$ or more against the
stock jaws' surviving conditions; all of these ratios are lower bounds,
because the finray ranges are right-censored at the analysis-window end.
For the application domains that motivated this work (laundry and
textile handling, hospital bedding, assistive robotics, garment
manufacturing, and soft-goods packaging), this means the flat-cloth
pick, the step that currently keeps such systems human-started, becomes
accessible to a standard arm running an unchanged policy, before any
investment in cameras, tactile sensing, or learned control. The
accompanying design rule is equally simple: press the compliant finger
past its buckling peak ($\approx 2$--$3$\,mm beyond first contact), and
the grasp becomes robust to both depth and tilt placement error.

\section{Future Work}

\begin{description}
  \item[Independent force-sensor cross-check.] All force measurements
        derive from the arm's joint-torque-based external-wrench
        estimate. Repeating a subset of the matrix with a dedicated
        force/torque sensor at the wrist (or an instrumented surface)
        would bound the estimator's pose-dependent bias directly.
  \item[Task-level evaluation.] The descoped downstream tasks
        (Chapter~\ref{ch:introduction}) --- folding from a corner grasp,
        placing onto curved objects, dragging between positions ---
        would test whether the contact-phase advantage measured here
        translates into end-to-end task success, including under learned
        policies.
  \item[Task-level precision comparison.] Beyond success rates, a soft
        versus rigid comparison of task-level \emph{precision} (placement
        accuracy after fold or drop) would characterise the cost side of
        compliance.
  \item[Dynamic and high-rate behaviour.] All trials here are
        quasi-static. The perturbation behaviour of compliant fingers
        under fast actuation and high-frequency control --- where the
        compliance that absorbs placement error may instead cause
        oscillation or delayed settling --- is uncharacterised.
  \item[Tilt-sensitivity model.] An analytical model relating finger
        curvature, contact compliance, and tilt-induced contact-point
        migration would connect the present empirical findings to the
        broader continuum-mechanics literature on soft contact.
  \item[Multi-fabric study.] All trials used a single cloth specimen.
        Replicating the matrix on fabrics of varying thickness, weave,
        and surface treatment would establish the applicability range
        of the depth--tilt-robustness curves identified here.
  \item[Closed-loop control.] The boundary-depth regime, where finray
        directional asymmetry dominates, motivates a tilt-compensating
        controller using inline force feedback, validated against the
        open-loop boundary baseline.
  \item[Sim2real validation.] A simulation of the finray finger in a
        compliant-contact physics engine could be calibrated against the
        empirical force--depth--tilt surface from this thesis, providing
        a digital twin for the gripper variants.
\end{description}
